\documentclass[10pt,letterpaper]{article}
\usepackage{cornell-notes}

\title{Chapter 6: Deadlocks}
\author{}
\date{}

\setDocTitle{Chapter 6: Deadlocks}
\setDocSubtitle{Operating Systems}
\setDocVersion{v1.0}
\setDocStatus{Study Companion}
\setDocOwner{Operating Systems Cornell Notes}
\setDocAuthor{}
\setDocDate{\today}
\setDocSummary{Cornell-format study and review notes for Chapter 6: Deadlocks.}

\setCornellCollection{Operating Systems Cornell Notes}
\setCornellUnitType{Chapter}
\setCornellUnitNumber{6}
\setCornellUnitTitle{Deadlocks}
\setCornellCompanionLabel{Study and review companion}

\hypersetup{pdftitle={Chapter 6: Deadlocks},pdfsubject={Cornell notes},pdfauthor={}}

\begin{document}
\maketitle
\section*{Learning Objectives}
\begin{studybox}{By the end of this chapter, you should be able to}
\begin{enumerate}
  \item Distinguish preemptable from nonpreemptable resources and safe acquisition protocols.
  \item State the four necessary conditions for resource deadlock.
  \item Use resource-allocation graphs to identify cycles and possible deadlock.
  \item Compare deadlock detection for single-instance and multiple-instance resource types.
  \item Apply safe-state reasoning and the single- or multiple-resource banker's algorithm.
  \item Explain how prevention attacks each necessary condition.
  \item Distinguish deadlock, communication deadlock, livelock, and starvation.
\end{enumerate}
\end{studybox}

\begin{studybox}{Section Map}
\hyperlink{sec-6-1}{6.1} \quad \hyperlink{sec-6-2}{6.2} \quad \hyperlink{sec-6-3}{6.3} \quad \hyperlink{sec-6-4}{6.4} \quad \hyperlink{sec-6-5}{6.5} \quad \hyperlink{sec-6-6}{6.6} \quad \hyperlink{sec-6-7}{6.7} \quad \hyperlink{sec-6-8}{6.8} \quad \hyperlink{sec-6-9}{6.9}
\end{studybox}

\section*{Cornell Cue-and-Notes Area}
\small
\begin{longtable}{C N}
\rowcolor{navy}
\color{white}\textbf{Cue / Recall Prompt} &
\color{white}\textbf{Notes}\\
\endfirsthead
\rowcolor{navy}
\color{white}\textbf{Cue / Recall Prompt} &
\color{white}\textbf{Notes (continued)}\\
\endhead
\rowcolor{ice}\multicolumn{2}{l}{\hypertarget{sec-6-1}{}\textbf{Section 6.1: Resources}}\\*\hline
\textbf{6.1.1 Resource classes} & A preemptable resource can be taken and restored without harm; a nonpreemptable resource cannot safely be revoked before use completes. Deadlock usually concerns exclusive nonpreemptable resources.\\\hline
\textbf{6.1.2 Acquisition} & A process requests, uses, and releases resources. If a request cannot be granted, it waits; correctness depends on acquisition order, partial holdings, and whether requests are known in advance.\\\hline
\rowcolor{ice}\multicolumn{2}{l}{\hypertarget{sec-6-2}{}\textbf{Section 6.2: Introduction to Deadlocks}}\\*\hline
\textbf{Deadlock definition} & A set is deadlocked when every member waits for an event that only another member of the same set can cause, so none can make progress.\\\hline
\textbf{6.2.1 Necessary conditions} & Resource deadlock requires mutual exclusion, hold and wait, no preemption, and circular wait at the same time. Removing any one condition prevents this form of deadlock.\\\hline
\textbf{6.2.2 Graph model} & Resource-allocation graphs use process and resource nodes with request and assignment edges. With one instance per resource type, a cycle is sufficient for deadlock; with multiple instances it is not.\\\hline
\rowcolor{ice}\multicolumn{2}{l}{\hypertarget{sec-6-3}{}\textbf{Section 6.3: The Ostrich Algorithm}}\\*\hline
\textbf{Deliberate tolerance} & A system may ignore rare deadlocks when prevention or detection costs more than occasional recovery by users or administrators. This is an economic policy, not a correctness proof.\\\hline
\rowcolor{ice}\multicolumn{2}{l}{\hypertarget{sec-6-4}{}\textbf{Section 6.4: Deadlock Detection and Recovery}}\\*\hline
\textbf{6.4.1 One resource each} & Build a wait-for or allocation graph and search for cycles. A cycle identifies the deadlocked set when every resource type has one instance.\\\hline
\textbf{6.4.2 Multiple resources} & Maintain available, current allocation, and outstanding request vectors. Repeatedly mark any process whose remaining request can be satisfied; unmarked processes after convergence are deadlocked.\\\hline
\textbf{6.4.3 Recovery} & Recovery can preempt resources, roll back to checkpoints, or terminate selected processes. Victim choice weighs priority, work lost, resources held, and repeated victimization.\\\hline
\rowcolor{ice}\multicolumn{2}{l}{\hypertarget{sec-6-5}{}\textbf{Section 6.5: Deadlock Avoidance}}\\*\hline
\textbf{Avoidance premise} & Avoidance grants only requests that leave a safe state, requiring advance knowledge of maximum claims or an equivalent future-demand bound.\\\hline
\textbf{6.5.1 Resource trajectories} & A two-process trajectory plots cumulative resource progress. Forbidden regions represent simultaneous claims, and a path must avoid entering a position from which completion cannot be guaranteed.\\\hline
\textbf{6.5.2 Safe states} & A state is safe if some completion order lets every process obtain its remaining maximum and finish. Unsafe does not mean currently deadlocked; it means no guarantee remains.\\\hline
\textbf{6.5.3 Single-resource banker} & Tentatively grant a request, then test whether available units can satisfy some process's remaining claim, reclaim its allocation, and repeat until all can finish.\\\hline
\textbf{6.5.4 Multiple-resource banker} & Generalize the safety test to vectors. A request must not exceed need or availability, and the tentative allocation is committed only if a full safe sequence still exists.\\\hline
\rowcolor{ice}\multicolumn{2}{l}{\hypertarget{sec-6-6}{}\textbf{Section 6.6: Deadlock Prevention}}\\*\hline
\textbf{6.6.1 Mutual exclusion} & Make resources sharable where semantics permit, for example by spooling output. Intrinsically exclusive resources cannot use this attack.\\\hline
\textbf{6.6.2 Hold and wait} & Require all resources at once or require release before requesting more. This prevents partial holdings but reduces utilization and may cause long waits.\\\hline
\textbf{6.6.3 No preemption} & If a new request cannot be granted, force release or rollback of held resources where their state can be restored safely.\\\hline
\textbf{6.6.4 Circular wait} & Assign a total order to resource types and require acquisition in increasing order. A cycle would require an impossible decrease back to the first resource.\\\hline
\rowcolor{ice}\multicolumn{2}{l}{\hypertarget{sec-6-7}{}\textbf{Section 6.7: Other Issues}}\\*\hline
\textbf{6.7.1 Two-phase locking} & A transaction first acquires all needed locks and then releases them without acquiring new ones. The protocol supports serializable execution but transactions can still deadlock during the growing phase.\\\hline
\textbf{6.7.2 Communication deadlocks} & Processes may wait cyclically for messages rather than reusable resources. Timeouts, protocol design, buffering, or cancellation can break waits, though timeouts can confuse delay with failure.\\\hline
\textbf{6.7.3 Livelock} & Processes remain active and respond to one another but repeatedly change state without useful progress. Randomized or asymmetric backoff can break the symmetry.\\\hline
\textbf{6.7.4 Starvation} & A process waits indefinitely because policy repeatedly favors others even though the system progresses. Fair queues or aging can bound postponement.\\\hline
\rowcolor{ice}\multicolumn{2}{l}{\hypertarget{sec-6-8}{}\textbf{Section 6.8: Research on Deadlocks}}\\*\hline
\textbf{Research themes} & Research studies distributed detection, database and transactional conflicts, static lock-order analysis, model checking, recovery policies, and deadlock behavior in communication protocols.\\\hline
\rowcolor{ice}\multicolumn{2}{l}{\hypertarget{sec-6-9}{}\textbf{Section 6.9: Summary}}\\*\hline
\textbf{Four strategy families} & Systems may ignore deadlock, detect and recover, avoid unsafe states, or prevent one necessary condition. Communication deadlock, livelock, and starvation require related but distinct reasoning.\\\hline
\end{longtable}
\normalsize

\section*{Chapter Synthesis}
\begin{studybox}{Chapter Summary}
Deadlock is permanent circular waiting within a set of processes. Resource-allocation models expose the conditions and cycles that cause it. A system can tolerate it, detect and recover, avoid unsafe allocations, or prevent a necessary condition; each strategy has information, runtime, utilization, and recovery costs.
\end{studybox}

\begin{studybox}{Key Relationships}
\begin{itemize}
  \item Resource requests create wait edges; allocations create ownership edges; cycles connect ownership to waiting.
  \item Detection reasons about the current state, whereas avoidance also considers declared future maximum demand.
  \item Prevention changes resource semantics or acquisition rules so a necessary condition cannot arise.
  \item Fairness mechanisms that combat starvation can also influence recovery-victim selection and retry behavior.
\end{itemize}
\end{studybox}

\begin{studybox}{Design Tradeoffs}
\begin{itemize}
  \item Ignoring deadlock minimizes routine overhead but accepts unbounded recovery disruption.
  \item Avoidance preserves safety but needs maximum claims and may deny requests that would have completed safely in practice.
  \item All-at-once acquisition prevents hold-and-wait but lowers utilization and concurrency.
  \item Frequent detection reduces deadlock duration but consumes more CPU and bookkeeping effort.
\end{itemize}
\end{studybox}

\begin{studybox}{Common Confusions}
\begin{itemize}
  \item An unsafe state is not necessarily deadlocked; it lacks a guaranteed completion sequence.
  \item A graph cycle proves deadlock only in the single-instance resource model.
  \item Deadlocked processes are blocked, livelocked processes remain active, and a starved process is bypassed while others progress.
\end{itemize}
\end{studybox}

\section*{Key Terms}
\small
\begin{longtable}{C N}
\rowcolor{navy}\color{white}\textbf{Term} & \color{white}\textbf{Concise definition}\\
\endfirsthead
\rowcolor{navy}\color{white}\textbf{Term} & \color{white}\textbf{Concise definition (continued)}\\
\endhead
\textbf{Deadlock} & Permanent circular waiting among a set of processes.\\\hline
\textbf{Preemptable resource} & Resource that can be revoked and later restored safely.\\\hline
\textbf{Nonpreemptable resource} & Resource that cannot be revoked without failure or lost work.\\\hline
\textbf{Mutual exclusion} & Condition in which only one process may use a resource at a time.\\\hline
\textbf{Hold and wait} & Holding resources while waiting to acquire additional resources.\\\hline
\textbf{Circular wait} & Cycle in which each process waits for a resource held by the next.\\\hline
\textbf{Safe state} & State with at least one guaranteed completion sequence.\\\hline
\textbf{Safe sequence} & Order in which every process can obtain remaining claims and finish.\\\hline
\textbf{Banker's algorithm} & Avoidance algorithm that grants requests only when the state remains safe.\\\hline
\textbf{Wait-for graph} & Process-only graph whose edges represent waiting on another process.\\\hline
\textbf{Rollback} & Restoration to an earlier checkpoint to release resources and retry.\\\hline
\textbf{Livelock} & Continued activity without useful forward progress.\\\hline
\textbf{Starvation} & Indefinite denial caused by allocation policy.\\\hline
\textbf{Two-phase locking} & Lock protocol with a growing acquisition phase and shrinking release phase.\\\hline
\end{longtable}
\normalsize

\enlargethispage{2\baselineskip}
\begin{studybox}{Active-Recall Review}
\small
\begin{enumerate}
  \item Which four conditions are necessary for resource deadlock?
  \item When does a cycle in a resource-allocation graph prove deadlock?
  \item How does multiple-instance detection identify processes that can finish?
  \item What factors influence recovery-victim choice?
  \item Why can an unsafe state still execute without deadlock?
  \item What extra information does the banker's algorithm require?
  \item How does resource ordering eliminate circular wait?
  \item How do communication deadlock, livelock, and starvation differ?
\end{enumerate}
\end{studybox}

\par\medskip
{\color{navy}\bfseries Next-Chapter Connections}\par\smallskip
Chapter 7 uses controlled resource multiplexing and isolation to run complete virtual machines and to support migration, checkpointing, and cloud services.
\normalsize
\end{document}
